\documentclass[conference]{IEEEtran}
\usepackage{amsmath,amssymb,graphicx,booktabs,url}
\usepackage[hidelinks]{hyperref}
\title{Short-Horizon PM2.5 Forecasting with Gradient Boosting:\\A Chronological Machine Learning Evaluation}
\author{\IEEEauthorblockN{Project 19 Team}
\IEEEauthorblockA{CSAI Student Project 19 --- Phase 1}}
\newcommand{\resultroot}{../../artifacts/phase1}
\IfFileExists{configs/phase1.toml}{\renewcommand{\resultroot}{artifacts/phase1}}{}
\begin{document}
\maketitle
\begin{abstract}
We study one-hour-ahead PM2.5 forecasting on the UCI Beijing hourly dataset using XGBoost with historical pollution, trailing summaries, calendar values, and weather. Initial training uses 2010--2012, parameter selection uses four expanding 2013 quarters, and the frozen settings are evaluated in four 2014 backtests. On 8,531 common observed target hours, XGBoost achieves MAE 11.833 and RMSE 21.088 $\mu$g/m$^3$, with $R^2=0.9489$. Its RMSE is 4.29\% lower than persistence. Ridge is slightly better overall, at RMSE 21.063, while daily seasonal naive performs substantially worse. A validation-only feature comparison shows incremental gains from temporal summaries and weather within the selected XGBoost configuration. However, XGBoost has higher error than persistence on high-pollution hours. These results support a modest short-horizon improvement while identifying peak prediction as a remaining weakness.
\end{abstract}
\begin{IEEEkeywords}
PM2.5, air quality forecasting, XGBoost, time series, chronological validation
\end{IEEEkeywords}

\section{Problem and Scope}
Our goal is to estimate $y_{t+h}$ at a fixed horizon of $h=1$ hour using measurements available through forecast origin $t$. The target is concentration in $\mu$g/m$^3$. We assume the hourly readings have been published before issuance; the archive does not provide the publication delays needed to verify this assumption in a live setting. This phase focuses on Machine Learning.

\section{Related Work}
Liang et al. study Beijing pollution and meteorological effects~\cite{liang2015}, giving us a reason to examine weather and seasonality. Their analysis does not supply a directly comparable one-hour forecasting score. Chen and Guestrin explain regularized tree boosting and sparse-data handling~\cite{chen2016}, which underpin our nonlinear model. Hyndman and Athanasopoulos~\cite{hyndman2021} guide our reference forecasts and chronological evaluation. Ridge regression~\cite{hoerl1970} provides a linear comparison suited to correlated lag features. Together, these sources shape our question: what do environmental features add beyond recent pollution and daily repetition?

\section{Dataset and Exploratory Analysis}
The selected source is UCI Beijing PM2.5, dataset 381~\cite{uci381}. The supplied CSV has 43,824 hourly records covering 2010--2014. PM2.5 was recorded at the US Embassy in Beijing, while weather was measured at Beijing Capital International Airport. Covariates include dew point, temperature, pressure, wind direction, accumulated wind speed, and snow/rain duration. This spatial mismatch is a limitation.

The timestamps are sorted and unique, with no absent hourly rows. PM2.5 has 2,067 missing measurements (4.72\%); the weather fields are complete. These gaps form 214 runs, with a maximum length of 155 hours. The loader validates dates and numeric values and preserves the regular grid. Missing observations remain explicit, and the ID is excluded. Every run records a CSV checksum and provenance. We do not interpolate evaluation targets or automatically remove extreme pollution readings.

EDA uses only the 26,304 initial training hours. Of these, 24,418 have measured PM2.5 and 1,886 are missing (7.17\%). The mean is 97.82, the median is 73.00, and the observed range is 0--994 $\mu$g/m$^3$. Skewness is 1.70; 797 readings exceed the upper IQR fence of 305.50. We retain these values because a large concentration alone does not establish a sensor error. The training 90th percentile, 218, later defines the high-pollution diagnostic.

The pooled hourly average is highest at 01:00, and October has the highest monthly average across the represented training years. One-hour autocorrelation is 0.961, falling to 0.394 at 24 hours and approximately $-0.006$ at 168 hours. Persistence is therefore a strong reference, while a weekly dependency is not established by this statistic. Accumulated wind speed has the largest absolute numeric weather correlation, $r=-0.264$; dew point has $r=0.245$. These are associations rather than causal effects. Fig.~\ref{fig:eda} shows the training concentration distribution.

\begin{figure}[ht]
\centering
\includegraphics[width=\columnwidth]{\resultroot/figures/article_distribution.pdf}
\caption{Training-period concentrations and the 90th-percentile threshold. Valid peaks are retained; missing observations are not interpolated.}
\label{fig:eda}
\end{figure}

\section{Feature Engineering}
Rows are indexed by forecast origin. Concentration lags include offsets $0,1,2,3,6,12,24,48,72,168$ hours. Lag zero is valid because the label is a future observation. For windows $w\in\{3,6,12,24,72,168\}$, compute means, population standard deviations and observed fractions over $[t-w+1,t]$. The observed fraction exposes reduced support when a window contains gaps. No future-based interpolation is performed.

Known target hour, weekday and month use sine/cosine encodings, with a weekend flag. Numeric weather is supplied at offsets 0, 1 and 24; missingness and categorical wind indicators are retained. This produces 70 columns before wind encoding and 73 after encoding the four observed wind categories. Learned encoders are fitted separately in each training fold. Ridge additionally uses fold-fitted median imputation and standardization. XGBoost uses native numeric missing-value handling. The first 168 elapsed hours are a warm-up.

\section{Models and Theory}
Persistence predicts $\hat y_{t+h|t}=y_t$. Seasonal naive with period $m=24$ predicts
\begin{equation}
\hat y_{t+h|t}=y_{t+h-m\lceil h/m\rceil}.
\end{equation}
At $h=1$, the seasonal source is $t-23$, exactly one day before the target's hour. Ridge minimizes squared residuals plus $\alpha\lVert\beta\rVert_2^2$, with fixed $\alpha=10$.

Gradient boosting constructs an additive ensemble $F_K(x)=b+\sum_{k=1}^{K}\eta f_k(x)$. XGBoost penalizes tree complexity and leaf weights:
\begin{equation}
\mathcal L=\sum_i \ell(y_i,F_K(x_i))+\sum_k\left(\gamma T_k+\frac{\lambda}{2}\lVert w_k\rVert_2^2\right).
\end{equation}
A second-order expansion uses gradient and Hessian sums $G,H$, giving regularized leaf weight $-G/(H+\lambda)$~\cite{chen2016}. Tree depth, learning rate, minimum child weight, row/column sampling and regularization control complexity. A fixed three-candidate grid selects parameters by pooled validation RMSE. The selected candidate has 300 trees, depth 3, learning rate 0.05, minimum child weight 5, row sampling 0.8, column sampling 0.9, and $\lambda=5$. Histogram trees use seed 42 and two threads. There is no test-based early stopping. All learned-model outputs are clipped at zero consistently.

\section{Evaluation Protocol}
Initial training uses 2010--2012. Four expanding quarterly folds in 2013 select parameters. Four quarterly 2014 folds evaluate the frozen setting. For a block beginning at target time $B$, the first origin is $B-h$. Training labels must satisfy $t_{\mathrm{label}}\le B-h$. This prevents future labels entering a fit, including at multi-hour horizons.

Models are refitted once per quarter; observed histories update hourly. Later folds may use earlier observed 2014 records according to this rolling policy. Targets and baseline inputs must be observed for a row to enter the common comparison. No evaluation target is imputed. Coverage and overlapping exclusion counts accompany scores.

For errors $e_i=\hat y_i-y_i$, the metrics are
\begin{align}
\mathrm{MAE}&=\frac1n\sum_i |e_i|,\\
\mathrm{RMSE}&=\sqrt{\frac1n\sum_i e_i^2},\\
R^2&=1-\frac{\sum_i e_i^2}{\sum_i(y_i-\bar y)^2}.
\end{align}
$R^2$ can be negative and is undefined for constant truth. Pooled metrics weight observed forecast hours equally; per-fold metrics expose temporal variability. Model selection and final evaluation are separate. Feature ablations use 2013 only with the selected settings held fixed, without separately optimizing each feature group.

\section{Results and Diagnostics}
\subsection{Validation and Feature Comparison}
We evaluated three XGBoost settings on the same 8,563 validation hours in 2013. Their pooled RMSEs were 26.772, 26.801 and 26.953. The selected model still trails persistence (24.870) and Ridge (23.809). In particular, XGBoost's first validation quarter has RMSE 42.471, versus 35.762 for persistence. The remaining quarters are easier for this model. We retain the original grid and report this weakness rather than altering the protocol after viewing the holdout.

Table~\ref{tab:ablation} shows that trailing/calendar features and then weather each improve the selected XGBoost configuration. This supports their incremental predictive value in this setup; it does not establish a causal effect or optimality for each feature set.
\begin{table}[ht]
\centering
\caption{Feature comparison on 2013 validation dates only.}
\label{tab:ablation}
\resizebox{\columnwidth}{!}{\input{\resultroot/report_tables/ablation.tex}}
\end{table}

\subsection{Rolling 2014 Backtest}
The common evaluation contains 8,531 of 8,760 expected target hours (97.39\% coverage). The other 229 hours lack an observed target or at least one baseline input. Counts for individual exclusions overlap.

\begin{table}[ht]
\centering
\caption{Measured rolling holdout errors. MAE/RMSE are in $\mu$g/m$^3$.}
\label{tab:overall}
\input{\resultroot/results_table.tex}
\end{table}

XGBoost reduces pooled RMSE by 4.29\% relative to persistence, while Ridge reduces it by 4.41\% (Table~\ref{tab:overall}). Ridge's advantage over XGBoost is only 0.025 $\mu$g/m$^3$. We have not tested the statistical significance of this small difference. Daily seasonal naive has $R^2=-0.1205$, indicating that simply repeating the previous day's target-hour value is weak despite the average calendar patterns.

Both learned models improve RMSE over persistence in every 2014 quarter, but their ordering changes (Table~\ref{tab:quarterly}). XGBoost is better than Ridge in Q2 and Q4; Ridge is better in Q1 and Q3. In Q1, XGBoost's MAE is worse than persistence despite its lower RMSE. Improvement therefore depends on both the period and the error measure.
\begin{table}[ht]
\centering
\caption{2014 RMSE by target quarter on common hours.}
\label{tab:quarterly}
\resizebox{\columnwidth}{!}{\input{\resultroot/report_tables/quarterly_rmse.tex}}
\end{table}

\subsection{High-Pollution Errors and Residuals}
There are 865 common holdout hours at or above 218 $\mu$g/m$^3$, the threshold fixed from initial training. XGBoost's RMSE on this subset is 39.959, compared with 38.557 for persistence and 38.522 for Ridge. Its bias is $-11.912$, while Ridge's is $-13.313$. Both learned models underpredict these high observations on average. In contrast, XGBoost has the lowest RMSE below the threshold (17.740). Much of its overall gain therefore comes from the lower-pollution subset.

For example, the recorded PM2.5 rises from 80 to 580 between 19:00 and 20:00 on 9 April 2014. XGBoost predicts 86.801 and Ridge predicts 92.469. This is XGBoost's largest holdout miss. Without an independent quality flag, we retain the observation rather than correcting it after inspecting the errors.

\begin{figure}[ht]\centering
\includegraphics[width=\columnwidth]{\resultroot/figures/article_peak_errors.pdf}
\caption{Errors below and at/above the training p90 threshold. XGBoost's pooled improvement does not extend to the high-pollution subset.}\end{figure}

One-hour residual autocorrelation is 0.120 for XGBoost and 0.063 for Ridge, with 24-hour values of 0.027 and 0.013. Some short-range structure remains. The final deployment XGBoost refit assigns 51.1\% of normalized gain importance to the latest concentration and 24.5\% to the three-hour mean. These importances describe the final refit, not causal effects or stability across quarterly fits. Additional outputs retain errors by month/hour, full timestamped predictions and a final-week forecast plot.

\section{Reproducibility}
The modular Python pipeline records data hash, source hash, configuration, versions, parameter decisions, fold audits, predictions and metrics. The supplied CSV's SHA256 begins \texttt{4127f868775e31b3}; the full digest is saved with the experiment. The run uses Python 3.13.5, pandas 2.2.3, scikit-learn 1.6.1 and XGBoost 3.4.1. A completed-run manifest is written only after the experiment outputs are produced. Deployment models are refitted after evaluation and are distinct from models used for backtest predictions. Tests cover feature causality, horizon alignment, missingness, fold boundaries, preprocessing and artifact generation. Cleaning exports and a worked feature preview make the data transformations inspectable.

\section{Limitations and Conclusion}
Historical performance at one Beijing location does not establish generalization to present-day conditions or another city. Missing baseline inputs restrict the comparison sample. Measured weather availability is assumed rather than established from an operational publication log. The finite grid does not guarantee optimal parameters. No calibrated uncertainty, causal effect, or valid daily AQI category is established.\par
The measured backtests show a modest pooled one-hour forecasting improvement over persistence for both XGBoost and Ridge, with Ridge slightly better overall. Temporal summaries and weather improve the selected XGBoost configuration on validation dates, but do not guarantee superiority over simpler models. High-pollution errors remain a clear weakness. The holdout has now been inspected; further changes motivated by these findings require a fresh evaluation or an explicit exploratory interpretation. This Phase 1 establishes the classical ML comparison without making deep-learning or hybrid-performance claims.

\bibliographystyle{IEEEtran}
\IfFileExists{reports/ieee/references.bib}{\bibliography{reports/ieee/references}}{\bibliography{references}}
\end{document}
